\documentclass{article}
\usepackage[utf8]{inputenc}
\usepackage[T1]{fontenc}
\usepackage{amsmath, amssymb, amsthm}
\usepackage{geometry}
\usepackage{hyperref}
\usepackage{bookmark}
\usepackage{tikz}


% \renewcommand{\thesubsection}{\Roman{subsection}}
% \newcommand{\exo}[2]{\textbf{\Roman{#1} #2} }

% \geometry{left=2.5cm, right=2.5cm, top=2.5cm, bottom=2.5cm}


\setlength{\parindent}{0pt}
\setlength{\parskip}{0.8em}

\title{Problèmes de Cauchy}
\author{Axel}
\date{\today}


\newcommand{\exo}[1]{
    \vspace{5mm}
    \textbf{Exercice #1}
    \hrulefill
}

\newcommand{\smalltext}[1]{
    \scriptsize
    \textit{#1}
    \normalsize
}

\begin{document}

\maketitle
% \tableofcontents

\section{Introduction}
dans ce document, nous allons étudier les problèmes de Cauchy.
Nous allons commencer par définir ce qu'est un problème de Cauchy,
la methode de résolution et enfin nous allons résoudre quelques exemples.

\subsection{Définition}
Un problème de Cauchy est un problème de résolution d'une équation différentielle
avec une condition initiale. En d'autres termes, il s'agit de trouver une fonction
qui vérifie une équation différentielle et qui vérifie une condition initiale.

\textit{exemple 1}
\begin{equation}
    \begin{cases}
        y'(x) = 2x \\
        y(0) = 1
    \end{cases}
\end{equation}

pour résoudre les problèmes de Cauchy, il est impératif de savoir primitiver une fonction.
dans l'exemple 1, la fonction $y(x)$ est la primitive de la fonction $2x$.

\section{Méthode de résolution}

\subsection{ordre 1}

les problèmes de Cauchy d'ordre 1 sont de la forme suivante

\begin{equation}\tag{E}\label{E}
    \begin{cases}
        y'(x) + a(x)y(x) = b(x) \\
        y(x_0) = y_0
    \end{cases}
\end{equation}

pour résoudre ce type de problème, il faut suivre les étapes suivantes

\begin{enumerate}
    \item trouver la solution homogene de l'équation différentielle
    \item trouver la solution particulière de l'équation différentielle
    \item trouver la solution du problème de Cauchy
\end{enumerate}

l'équation homogène associée à l'équation différentielle \eqref{E} est:

\begin{equation}\tag{H}
    y'(x) + a(x)y(x) = 0
\end{equation}

les solutions de l'équation homogène sont de la forme $y_h(x) = Ce^{-A(x)}$ où $C$ est une constante et $A(x)$ est une primitive de $a(x)$.

Pour trouver la solution particulière de l'équation différentielle \eqref{E}, il faut trouver une fonction $y_p(x)$ qui vérifie l'équation différentielle \eqref{E}.

en utilisant la méthode de variation de la constante,
la solution particulière de l'équation différentielle \eqref{E} est de la forme 
\begin{equation*}
    y_p(x) = u(x)y_h(x)
\end{equation*}
où $u(x)$ est une primitive de $\frac{b(x)}{y_h(x)}$.

Ainsi, la solution générale de l'équation différentielle \eqref{E} est de la forme
\begin{equation*}
    y(x) = y_h(x) + y_p(x)
\end{equation*}

il ne reste plus qu'à trouver la constante $C$ en utilisant la condition initiale $y(x_0) = y_0$.

\subsection{ordre 2}
Les problèmes de Cauchy d'ordre 2 sont de la forme suivante :
\begin{equation}
    y''(x) + a(x)y'(x) + b(x)y(x) = c(x)
\end{equation}

Pour résoudre ce type de problème, suivez les étapes suivantes :
\begin{enumerate}
    \item Trouver les solutions de l'équation homogène associée.
    \item Trouver une solution particulière de l'équation non homogène.
    \item Combiner les solutions pour obtenir la solution générale.
    \item Utiliser les conditions initiales pour déterminer les constantes.
\end{enumerate}

L'équation homogène associée est :
\begin{equation}
    y''(x) + a(x)y'(x) + b(x)y(x) = 0
\end{equation}

Les solutions de l'équation homogène dépendent des racines de l'équation caractéristique :
\begin{equation}
    r^2 + ar + b = 0
\end{equation}

\begin{itemize}
    \item Si les racines sont réelles et distinctes, la solution est de la forme $y_h(x) = C_1 e^{r_1 x} + C_2 e^{r_2 x}$.
    \item Si les racines sont réelles et égales, la solution est de la forme $y_h(x) = (C_1 + C_2 x)e^{r x}$.
    \item Si les racines sont complexes, la solution est de la forme $y_h(x) = e^{\alpha x}(C_1 \cos(\beta x) + C_2 \sin(\beta x))$.
\end{itemize}

Pour trouver une solution particulière $y_p(x)$, utilisez des méthodes comme l'ansatz ou la variation des constantes.

La solution générale est :
\begin{equation}
    y(x) = y_h(x) + y_p(x)
\end{equation}

Utilisez les conditions initiales pour déterminer les constantes $C_1$ et $C_2$.



\subsection*{Shema}







\section{Exercices}


\exo{1}

\begin{equation*}
\begin{cases}
2y' + y = -4x+1\\
y(1)=0 
\end{cases}
\end{equation*}

\scriptsize
\textit{indication : nous pouvons chercher une solution particulière de la forme $y_p(x) = ax+b$}
\normalsize


\exo{2}

\begin{equation*}
\begin{cases}
y' + y = xe^{-x}\\
y(1)=\frac{1}{2}
\end{cases}
\end{equation*}

\scriptsize
\textit{indication : nous pouvons chercher une solution particulière avec la méthode de variation de la constante}
\normalsize


\exo{3}

\begin{equation*}
\begin{cases}
    y'' + y' + y = 0  \\
    y(0)=0
\end{cases}
\end{equation*}

\scriptsize
\textit{indication : il nous faudra passer par l'équation caractéristique à deux racines complexes}
\normalsize

\exo{4}

\begin{equation*}
\begin{cases}
    y' + xy = x\sin(x^2) \\
    y(0)=0
\end{cases}
\end{equation*}

\exo{5}

\begin{equation*}
\begin{cases}
    y'' - y' - 2y = x^2 \\
    y(0)=1
\end{cases}
\end{equation*}

\exo{6}

\begin{equation*}
\begin{cases}
    y' - ln(x)y = x^x \\
    y(1) = 1
\end{cases}
\end{equation*}


\section{Correction}
\subsection{Exercice 1}
Résoudre le problème de Cauchy suivant :
\begin{equation*}
\begin{cases}
2y' + y = -4x+1\\
y(1)=0 
\end{cases}
\end{equation*}

On pose l'équation homogène suivante :
\begin{equation*}
\begin{aligned}
2y' + y &= 0 \\
y' + \frac{1}{2}y &= 0 
\end{aligned}
\end{equation*}
On retrouve la forme $y'+a(x)y = 0$. On applique la formule triviale :

\begin{equation*}
y = \lambda e^{-A(x)} \quad  \lambda\in \mathbb{R}
\end{equation*}

où $A(x)$ est la primitive de $a(x)$. 

Ainsi, nous avons comme solution de l'équation homogène $y_h$:
\begin{equation*}
y_{h} = \lambda e^{-x/2}\quad \lambda\in \mathbb{R}
\end{equation*}

Cherchons une solution particulière. On remarque que $-4x+1$ est un polynôme, on cherche alors une fonction du même type.
Soit $a,b \in \mathbb{R}$, on suppose la solution particulière de la forme $y_{p} = ax+b$.
On pose l'équation suivante:
\begin{equation*}
\begin{aligned}
 2y'_{p} + y_{p} &= -4x+1 \\
 2a + ax + b &= -4x + 1
\end{aligned}
\end{equation*}
Nous avons donc le système suivant :
\begin{equation*}
\begin{gathered}
\begin{cases}
2a+b = 1\\
a = -4
\end{cases} \\
\Leftrightarrow \\
\begin{cases}
b = 9\\
a = -4
\end{cases}
\end{gathered}
\end{equation*}
Ainsi,
\begin{equation*}
y_{p} = -4x+9
\end{equation*}

On a donc la solution générale à notre équation :
\begin{equation*}
\lambda e^{-x/2} -4x+9\quad \lambda\in \mathbb{R}
\end{equation*}
Il s'agit maintenant de déterminer la valeur de $\lambda$ telle que $y(1)=0$.
On pose alors : 
\begin{equation*}
\begin{aligned}
\lambda e^{-1/2} -4+9 &= 0 \\
\lambda e^{-1/2}  &= -5 \\
\lambda   &= -5e^{1/2} 
\end{aligned}
\end{equation*}

Nous avons donc la solution au problème de Cauchy suivante :
\begin{equation*}
\begin{gathered}
 -5e^{1/2}e^{-x/2} -4x+9 \\
 -5e^{-x\frac{1}{2} + \frac{1}{2}} -4x+9 \\
 -5e^{\frac{1-x}{2}} -4x+9 
\end{gathered}
\end{equation*}

\begin{equation*}
y(x) = -5e^{\frac{1-x}{2}} -4x+9 
\end{equation*}

\subsection{Exercice 2}

Résoudre le problème de Cauchy suivant :
\begin{equation*}
\begin{cases}
y' + y = xe^{-x}\\
y(1)=\frac{1}{2}
\end{cases}
\end{equation*}

On pose l'équation homogène suivante :
\begin{equation*}
y_{0}' + y_{0} = 0
\end{equation*}
Les solutions de l'équation homogène sont de la forme :
\begin{equation*}
y_{0} = \lambda e^{-x}
\end{equation*}

En utilisant la méthode de variation de la constante, on cherche $\lambda(x)$ tel que :
\begin{equation*}
\lambda'(x)y_{0} = xe^{-x}
\end{equation*}
Ainsi,
\begin{equation*}
\lambda'(x) = x \implies \lambda(x) = \frac{x^2}{2}
\end{equation*}

On a donc une solution particulière $y_{p}$ :
\begin{equation*}
y_{p} = \frac{x^2}{2}e^{-x}
\end{equation*}

La solution générale de l'équation différentielle est donc :
\begin{equation*}
y(x) = \frac{x^2}{2}e^{-x} + \lambda e^{-x}
\end{equation*}

Il reste à déterminer la constante $\lambda$ en utilisant la condition initiale $y(1) = \frac{1}{2}$ :
\begin{equation*}
\frac{1}{2} = \frac{1}{2}e^{-1} + \lambda e^{-1} \implies \frac{1}{2} = \frac{1}{2e} + \frac{\lambda}{e}
\end{equation*}
\begin{equation*}
\frac{1}{2} = \frac{1}{2e} + \frac{\lambda}{e} \implies \frac{1}{2} - \frac{1}{2e} = \frac{\lambda}{e}
\end{equation*}
\begin{equation*}
\lambda = \frac{1}{2}(e - 1)
\end{equation*}

La solution au problème de Cauchy est donc :
\begin{equation*}
y(x) = \frac{x^2}{2}e^{-x} + \frac{1}{2}(e - 1)e^{-x}
\end{equation*}

\subsection{Exercice 3}
Résoudre le problème de Cauchy suivant :
\begin{equation*}
\begin{cases}
    y'' + y' + y = 0  \\
    y(0)=0
\end{cases}
\end{equation*}

\textbf{Étape 1 : Établir l'équation caractéristique}

L'équation différentielle donnée est une équation linéaire à coefficients constants. Pour la résoudre, nous associons l'équation caractéristique :
\begin{equation*}
r^2 + r + 1 = 0
\end{equation*}

\textbf{Étape 2 : Résoudre l'équation caractéristique}

Calculons le discriminant de cette équation quadratique :
\begin{equation*}
\Delta = b^2 - 4ac = (1)^2 - 4 \times 1 \times 1 = 1 - 4 = -3
\end{equation*}

Comme le discriminant est négatif, les racines sont complexes conjuguées :
\begin{equation*}
r = \frac{-b \pm i\sqrt{|\Delta|}}{2a} = \frac{-1 \pm i\sqrt{3}}{2}
\end{equation*}
où \( i \) est l'unité imaginaire.

\textbf{Étape 3 : Écrire la solution générale de l'équation différentielle}

Lorsque les racines de l'équation caractéristique sont complexes conjuguées de la forme \( r = \alpha \pm i\beta \), la solution générale de l'équation différentielle est :
\begin{equation*}
y(x) = e^{\alpha x} \left( A\cos(\beta x) + B\sin(\beta x) \right)
\end{equation*}
où \( A \) et \( B \) sont des constantes réelles.

Dans notre cas, \( \alpha = -\dfrac{1}{2} \) et \( \beta = \dfrac{\sqrt{3}}{2} \). Donc, la solution générale est :
\begin{equation*}
y(x) = e^{-\frac{x}{2}} \left( A\cos\left( \frac{\sqrt{3}}{2} x \right) + B\sin\left( \frac{\sqrt{3}}{2} x \right) \right)
\end{equation*}

\textbf{Étape 4 : Appliquer la condition initiale}

Nous disposons de la condition initiale \( y(0) = 0 \). Calculons \( y(0) \) :
\begin{align*}
y(0) &= e^{0} \left( A\cos\left( 0 \right) + B\sin\left( 0 \right) \right) \\
&= 1 \left( A \times 1 + B \times 0 \right) \\
&= A
\end{align*}
Étant donné que \( y(0) = 0 \), nous en déduisons :
\begin{equation*}
A = 0
\end{equation*}

\textbf{Étape 5 : Écrire la solution particulière avec la constante déterminée}

La solution devient donc :
\begin{equation*}
y(x) = e^{-\frac{x}{2}} B\sin\left( \frac{\sqrt{3}}{2} x \right)
\end{equation*}

\textbf{Remarque :} Pour déterminer la constante \( B \), une seconde condition initiale serait nécessaire, par exemple la valeur de \( y'(0) \). Sans cette information supplémentaire, \( B \) reste indéterminée.

\subsection{Exercice 4}
Résoudre le problème de Cauchy suivant :
\begin{equation*}
\begin{cases}
    y' + xy = x\sin(x^2) \\
    y(0)=0
\end{cases}
\end{equation*}

\textbf{Étape 1 : Trouver la solution homogène}

L'équation différentielle donnée est une équation linéaire à coefficients constants. Pour la résoudre, nous associons l'équation homogène :
\begin{equation*}
y' + xy = 0
\end{equation*}

Les solutions de l'équation homogène sont de la forme :
\begin{align*}
    y_h &= \lambda e^{-\!\int\!x} \\
    &= \lambda e^{-\frac{1}{2}x^2}
\end{align*}


\textbf{Étape 2 : Trouver la solution particulière}

Nous utilisons la méthode de variation de la constante :

\begin{align*}
y_p &= \lambda(x) e^{ - \frac{1}{2} x^2 } \\
\lambda'(x) &= x\sin(x^2) e^{ \frac{1}{2} x^2 }
\end{align*}

Calculons \(\lambda(x)\) en intégrant :

\begin{equation*}
\lambda(x) = \int x\sin(x^2) e^{ \frac{1}{2} x^2 } dx
\end{equation*}

Effectuons le changement de variable :

\begin{equation*}
t = x^2 \quad \Rightarrow \quad dt = 2x\,dx \quad \Rightarrow \quad x\,dx = \frac{dt}{2}
\end{equation*}

L'intégrale devient :

\begin{equation*}
\lambda(x) = \frac{1}{2} \int \sin(t) e^{ \frac{1}{2} t } dt
\end{equation*}

Calculons l'intégrale \(I\) :

\begin{equation*}
I = \int \sin(t) e^{ \frac{1}{2} t } dt
\end{equation*}

Utilisons la méthode DI pour l'intégration par parties :

\begin{center}
\begin{tabular}{c|c}
\textbf{D} (Dérive) & \textbf{I} (Intègre) \\
\hline
\(\sin(t)\) & \(e^{ \frac{1}{2} t }\) \\
\(\cos(t)\) & \(\frac{2}{1} e^{ \frac{1}{2} t }\) \\
\(-\sin(t)\) & \(\frac{4}{1} e^{ \frac{1}{2} t }\) \\
\(\cos(t)\) & \(\frac{8}{1} e^{ \frac{1}{2} t }\) \\
\end{tabular}
\end{center}

En pratique, il est plus efficace d'utiliser la formule standard pour cette intégrale :

\begin{equation*}
I = \int e^{ a t } \sin( b t ) dt = \frac{ e^{ a t } }{ a^2 + b^2 } \left( a \sin( b t ) - b \cos( b t ) \right ) + C
\end{equation*}

Dans notre cas, \(a = \dfrac{1}{2}\) et \(b = 1\), donc :

\begin{equation*}
I = \frac{ e^{ \frac{1}{2} t } }{ \left( \dfrac{1}{2} \right)^2 + 1 } \left( \dfrac{1}{2} \sin( t ) - \cos( t ) \right ) + C
\end{equation*}

Calculons le dénominateur :

\begin{equation*}
\left( \dfrac{1}{2} \right)^2 + 1 = \dfrac{1}{4} + 1 = \dfrac{5}{4}
\end{equation*}

Simplifions l'expression :

\begin{equation*}
I = \frac{ 4 }{ 5 } e^{ \frac{1}{2} t } \left( \dfrac{1}{2} \sin( t ) - \cos( t ) \right ) + C
\end{equation*}

Ainsi, nous obtenons :

\begin{equation*}
\lambda(x) = \frac{1}{2} I = \frac{2}{5} e^{ \frac{1}{2} x^2 } \left( \dfrac{1}{2} \sin( x^2 ) - \cos( x^2 ) \right ) + C
\end{equation*}

La solution particulière est donc :

\begin{equation*}
y_p = \lambda(x) e^{ - \frac{1}{2} x^2 } = \frac{2}{5} \left( \dfrac{1}{2} \sin( x^2 ) - \cos( x^2 ) \right ) + C e^{ - \frac{1}{2} x^2 }
\end{equation*}

La solution générale est :

\begin{equation*}
y(x) = y_h + y_p = D e^{ - \frac{1}{2} x^2 } + \frac{2}{5} \left( \dfrac{1}{2} \sin( x^2 ) - \cos( x^2 ) \right )
\end{equation*}

En appliquant la condition initiale \(y(0) = 0\) :

\begin{align*}
0 &= D e^{ 0 } + \frac{2}{5} \left( \dfrac{1}{2} \sin( 0 ) - \cos( 0 ) \right ) \\
0 &= D - \frac{2}{5} (1) \\
D &= \frac{2}{5}
\end{align*}

La solution du problème de Cauchy est donc :

\begin{equation*}
y(x) = \frac{2}{5} e^{ - \frac{1}{2} x^2 } + \frac{2}{5} \left( \dfrac{1}{2} \sin( x^2 ) - \cos( x^2 ) \right )
\end{equation*}

Simplifions :

\begin{equation*}
y(x) = \frac{2}{5} \left( e^{ - \frac{1}{2} x^2 } + \dfrac{1}{2} \sin( x^2 ) - \cos( x^2 ) \right )
\end{equation*}


\subsection{Exercice 5}

Résoudre le problème de Cauchy suivant :
\begin{equation*}
\begin{cases}
    y'' - y' - 2y = x^2 \\
    y(0)=1
\end{cases}
\end{equation*}


\textbf{Étape 1 : Trouver la solution homogène}

L'équation différentielle donnée est une équation linéaire à coefficients constants.
On a donc l'equation caractéristique suivante :

\begin{equation*}
r^2 - r - 2 = 0
\end{equation*}

$\Delta = 1 + 8 = 9$ donc les racines sont $r_1 = 2$ et $r_2 = -1$.

La solution homogène est donc :
\begin{equation*}
    y_h = \lambda e^{2x} + \mu e^{-x}
\end{equation*}

chercheons une solution particulière de l'équation.
On suppose que la solution particulière est de la forme $y_p = ax^2 + bx + c$.

On a donc:
\begin{equation*}
\begin{aligned}
    y_p &= ax^2 + bx + c \\
    y_p' &= 2ax + b \\
    y_p'' &= 2a
\end{aligned}
\end{equation*}

En remplaçant dans l'équation différentielle, on obtient:
\begin{equation*}
\begin{aligned}
    2a - 2ax - b - 2ax - b - 2(ax^2 + bx + c) &= x^2 \\
    -2ax^2 - 2bx - 2c &= x^2
\end{aligned}
\end{equation*}

On a donc le système suivant:
\begin{equation*}
\begin{aligned}
    -2a &= 1 \\
    -2(b + a) &= 0 \\
    2a - b - 2c &= 0
\end{aligned}
\end{equation*}

On trouve alors $a = -\frac{1}{2}$, $b = \frac{1}{2}$ et $c = -\frac{3}{4}$.

La solution particulière est donc:
\begin{equation*}
    y_p = -\frac{1}{2}x^2 + \frac{1}{2}x - \frac{3}{4}
\end{equation*}

La solution générale est donc:
\begin{equation*}
    y = \lambda e^{2x} + \mu e^{-x} -\frac{1}{2}x^2 + \frac{1}{2}x - \frac{3}{4}
\end{equation*}

En utilisant la condition initiale $y(0) = 1$, nous pouvons eliminer $\mu$ et trouver $\lambda$.

On pose alors:
\begin{equation*}
\begin{aligned}
    1 &= \lambda + \mu - \frac{3}{4}   \\
    \lambda &= 1 + \frac{3}{4} - \mu \\
    \lambda &= \frac{7}{4} - \mu
\end{aligned}
\end{equation*}

En remplaçant dans l'équation générale, on obtient:
\begin{equation*}
    y = (\frac{7}{4} - \mu)e^{2x} + \mu e^{-x} -\frac{1}{2}x^2 + \frac{1}{2}x - \frac{3}{4}
\end{equation*}



\subsection{Exercice 6}

Résoudre le problème de Cauchy suivant :
\begin{equation*}
\begin{cases}
    y' - \ln(x)y = x^x \\
    y(1) = 1
\end{cases}
\end{equation*}

\textbf{Étape 1 : Trouver la solution homogène}

L'équation différentielle donnée est une équation linéaire dont le coefficient est une fonction de \( x \).

Pour la résoudre, nous associons l'équation homogène :

\begin{equation*}
y' - \ln(x) y = 0
\end{equation*}

Les solutions de l'équation homogène sont de la forme :

\begin{equation*}
y_h = \lambda e^{-A(x)}
\end{equation*}

où \( A(x) \) est la primitive de \( \ln(x) \).

il nous faut donc trouver la primitive de \( \ln(x) \) qui est \( x \ln(x) - x \).

La solution homogène est donc :

\begin{align*}
y_h &= \lambda e^{x \ln(x) - x} \\
    &= \lambda e^{-x} x^x
\end{align*}

\textbf{Étape 2 : Trouver la solution particulière}

Nous utilisons la méthode de variation de la constante :

\begin{align*}
y_p &= \lambda(x) e^{-x} x^x \\
\end{align*}

Calculons \( \lambda(x) \):

\begin{align*}
\lambda'(x)(e^{-x} x^x) &= x^x \\
\lambda'(x) &= \frac{x^x}{e^{-x} x^x} \\
\lambda'(x) &= \frac{1}{e^{-x}} \\
\lambda'(x) &= e^{x} \\
\lambda(x) &= e^{x}
\end{align*}

La solution particulière est donc :

\begin{align*}
y_p &= e^{x} e^{-x} x^x \\
    &= x^x
\end{align*}

La solution générale est :

\begin{align*}
y(x) &= y_h + y_p \\
    &= \lambda e^{-x} x^x + x^x \\
    &= x^x (\lambda e^{-x} + 1)
\end{align*}

\textbf{Étape 3 : Appliquer la condition initiale}

Nous disposons de la condition initiale \( y(1) = 1 \). Calculons \( y(1) \):

\begin{align*} 
1^1 (\lambda e^{-1} + 1) &= 1\\
    \lambda e^{-1} + 1 &= 1 \\
    \lambda e^{-1} &= 0 \\
    \lambda &= 0
\end{align*}

La solution du problème de Cauchy est donc :

\begin{equation*}
    y(x) = x^x
\end{equation*}





\end{document}


